\appendix
\section{Full hyperparameters, schedule derivation and reproducibility}
\label{sec:appendix}

\subsection{Hyperparameters by family}

The framework groups models into three \emph{families}, each with its own base
configuration and per-phase overrides. Only the convolutional family is inside
this paper's grid; the vision-transformer and language-model families are
implemented, configured and runnable in the same framework, and are listed here
for completeness, but \textbf{no result in this paper comes from them} and they
are outside the claims of Section~\ref{sec:results}.

\begin{table}[h]
\centering
\small
\caption{Convolutional family --- ResNet-34, ResNet-50, VGG-11, VGG-19. This is
the family the paper's grid is drawn from. ResNet-101 is configured identically
and was not run in this grid.}
\label{tab:hyp-cnn}
\begin{tabular}{@{}p{0.26\textwidth}p{0.28\textwidth}p{0.42\textwidth}@{}}
\toprule
\multicolumn{3}{@{}p{0.96\textwidth}@{}}{\textit{Shared across all phases}} \\
\midrule
Dataset & \multicolumn{2}{p{0.70\textwidth}@{}}{CIFAR-10, 10 classes, $32\times32$ inputs} \\
Split & \multicolumn{2}{p{0.70\textwidth}@{}}{45{,}000 train / 5{,}000 val / 10{,}000 held-out test} \\
Batch size & \multicolumn{2}{p{0.70\textwidth}@{}}{512 (87 optimizer steps per epoch)} \\
Precision & \multicolumn{2}{p{0.70\textwidth}@{}}{bf16 autocast, fp32 master weights} \\
Gradient clipping & \multicolumn{2}{p{0.70\textwidth}@{}}{global norm 10.0} \\
Data loader workers & \multicolumn{2}{p{0.70\textwidth}@{}}{24} \\
\midrule
\multicolumn{3}{@{}p{0.96\textwidth}@{}}{\textit{Phase-specific}} \\
\midrule
 & \textbf{Dense} & \textbf{I.P.\ / BaCP pruning phase} \\
Optimizer & SGD & SGD \\
Learning rate & 0.01 & 0.01 (I.P.); 0.1 (BaCP; VGG-11 and VGG-19: 0.05) \\
LR schedule & linear with warmup & constant \\
Epochs & 100, early stop patience 20 & 50 \\
Interleaved recovery epochs & --- & 0 (recovery is the schedule's tail) \\
Sparsity schedule & --- & cubic \citep{zhu2018prune}, ramp 80\% / hold 20\% \\
Mask interval $\Delta T$ & --- & 100 optimizer steps \\
Prunable task head & --- & yes \\
WANDA comparison group & --- & global \\
\midrule
\multicolumn{3}{@{}p{0.96\textwidth}@{}}{\textit{BaCP-only}} \\
\midrule
Temperature $\tau$ & \multicolumn{2}{p{0.70\textwidth}@{}}{0.15} \\
$(\lambda_1,\lambda_2,\lambda_3,\lambda_4)$ & \multicolumn{2}{p{0.70\textwidth}@{}}{$(0.25, 0.25, 0.25, 0.25)$ on PrC, FiC, SnC, CE} \\
Contrastive formulation & \multicolumn{2}{p{0.70\textwidth}@{}}{\texttt{legacy}: SupCon $+$ NT-Xent over the square $2B \times 2B$ matrix} \\
Projection head & \multicolumn{2}{p{0.70\textwidth}@{}}{per-model trainable $\mathrm{Linear}(d, 128)$, never pruned, discarded before fine-tune} \\
\midrule
\multicolumn{3}{@{}p{0.96\textwidth}@{}}{\textit{Fine-tune (both arms, identical)}} \\
\midrule
Optimizer / LR / epochs & \multicolumn{2}{p{0.70\textwidth}@{}}{AdamW \citep{loshchilov2019adamw}, $10^{-4}$, 25 epochs, mask frozen} \\
Model selection & \multicolumn{2}{p{0.70\textwidth}@{}}{validation accuracy on the 5{,}000-image split} \\
\bottomrule
\end{tabular}
\end{table}

\begin{table}[h]
\centering
\small
\caption{Families implemented in the framework but \textbf{outside this paper's
grid}. Listed for completeness; no reported result uses them.}
\label{tab:hyp-other}
\begin{tabular}{@{}p{0.22\textwidth}p{0.30\textwidth}p{0.44\textwidth}@{}}
\toprule
 & \textbf{Vision transformer} & \textbf{Language model} \\
 & ViT-Tiny, ViT-Small & DistilBERT-base, RoBERTa-base \\
\midrule
Dataset & CIFAR-10, $224\times224$ & SST-2 (GLUE), 2 classes \\
Batch size & 256 & 64 \\
Initialisation & Hugging Face hub weights & Hugging Face hub weights \\
Dense phase & SGD $0.01$, linear warmup, 100 epochs, patience 20 & AdamW $2\times10^{-5}$, linear warmup, 5 epochs \\
I.P.\ phase & AdamW $10^{-3}$, 50 epochs & AdamW $5\times10^{-5}$, 15 epochs \\
BaCP phase & SGD $0.01$, 50 epochs & AdamW, 15 epochs; DistilBERT $5\times10^{-5}$, RoBERTa $10^{-5}$ \\
Mask interval $\Delta T$ & 175 steps ($=1$ epoch) & 1052 steps ($=1$ epoch) \\
Fine-tune & AdamW $10^{-4}$, 25 epochs & AdamW $10^{-4}$, 10 epochs \\
Task head / embeddings & pruned (vision convention) & kept dense, the PLM convention of \citet{xu2022cap} \\
Evaluation caveat & --- & GLUE's test split is unlabelled, so the reported ``test'' accuracy would be the validation set \\
\bottomrule
\end{tabular}
\end{table}

\subsection{Derivation of the sparsity schedule}

Let $T_{\mathrm{tot}}$ be the total number of optimizer steps in the pruning
phase, $T = \alpha T_{\mathrm{tot}}$ the end of the ramp with $\alpha = 0.8$,
$s_f$ the target sparsity, and $\Delta T$ the mask-update interval. The scheduled
sparsity at step $t$ is Eq.~\eqref{eq:cubic},
$s_t = s_f\big(1 - (1 - t/T)^3\big)$ for $t \le T$ and $s_t = s_f$ thereafter, the
schedule of \citet{zhu2018prune} with $s_i = 0$ and $t_0 = 0$. The mask is
recomputed at steps $t \in \{\Delta T, 2\Delta T, \dots\}$ with $t \le T$, and
re-applied after every optimizer step throughout.

For the convolutional family this instantiates as follows. With 45{,}000 training
images at batch 512 there are 87 optimizer steps per epoch, so a 50-epoch pruning
phase gives $T_{\mathrm{tot}} = 4{,}350$ steps. The ramp ends at
$T = 0.8 \times 4{,}350 = 3{,}480$ steps, which is 40 epochs, and the mask is held
fixed for the remaining 870 steps, which is 10 epochs of recovery. At
$\Delta T = 100$ the mask is therefore recomputed 34 times over the ramp.

Two properties of the cubic form are worth making explicit at extreme sparsity.
The derivative $ds_t/dt = 3 s_f (1 - t/T)^2 / T$ is largest at $t = 0$ and decays
to zero at the end of the ramp, so most of the \emph{sparsity} is reached early
and the late updates move $s_t$ very little. But the quantity the network
experiences is the fraction of \emph{surviving} weights removed, which is
$\big(s_{t+\Delta T} - s_t\big) / \big(1 - s_t\big)$, and at $s_f = 0.999$ that
ratio grows across the ramp even as the absolute increments shrink: the final
updates remove a large share of a very small population. This asymmetry is the
reason the ramp is not allowed to run to the last step. The $\alpha = 0.8$ split
guarantees a fixed 20\% of the budget in which no further weights are removed and
the surviving ones can re-fit, following \citet{gale2019state}.

The choice $\Delta T = 100$ rather than the exact per-epoch count of 87 is
deliberate and conservative: \citet{zhu2018prune} report the pruning frequency as
having negligible effect over the range 100 to 1000 steps, and 87 would place the
protocol below the smallest interval they measured. Because mask updates are
driven by optimizer steps rather than epoch boundaries, the schedule is invariant
to changes in the number of batches per epoch.

\subsection{Reproducibility}

\paragraph{Single source of protocol truth.}
Every hyperparameter in Tables~\ref{tab:protocol}, \ref{tab:hyp-cnn} and
\ref{tab:hyp-other} is read from one file, \texttt{project/test\_notebooks/nb\_common.py},
whose \texttt{FAMILIES} dictionary holds a base configuration per model family
plus per-phase overrides for the dense, I.P.\ and BaCP phases. Both arms of every
cell are constructed from that dictionary by the same builder function, which is
what makes the matched-budget claim of Section~\ref{sec:experiments} mechanical
rather than a matter of care: the I.P.\ and BaCP configurations are derived from a
shared base, so a protocol value cannot drift between arms without changing both.
The BaCP-specific keys (temperature, $\lambda$ vector, contrastive formulation,
projection-head mode, fine-tune optimizer) are the only additions in the BaCP
override.

\paragraph{Provenance enforcement.}
Initialisation is verified before training rather than assumed. The model factory
records what the loaded weights actually are --- an ImageNet checkpoint, hub
weights, or random initialisation --- and the notebooks' preflight check raises and
halts when the recorded provenance is not the requested one. This exists because a
checkpoint load that fails soft is otherwise indistinguishable from a run that
trained from scratch, and a ``pretrained'' baseline that was never pretrained
would invalidate every cell measured against it.

\paragraph{Run identity and completion.}
Each run is addressed by a structured key of the form
\texttt{static.\{phase\}.\{model\}.\{dataset\}.\{sparsity\}.\{pruner\}.seed\{n\}}.
A cell counts as complete if and only if a result record carrying its key exists;
deleting the record is the re-run mechanism. Failures write records with an
explicit failed status, so a crashed run is never silently indistinguishable from
one that was never scheduled --- which matters for a grid reported at partial
completion.

\paragraph{Evaluation.}
All reported accuracies are top-1 on the full 10{,}000-image CIFAR-10 test set,
evaluated once per completed cell, with no partial final batch dropped. Sparsity
is reported as the fraction of prunable weights that are masked, where the
prunable scope excludes 1-D tensors, frozen tensors and the projection head, and
includes the classifier head. The scope is identical between the two arms of every
cell.

\subsection{Seed and variance reporting}

The run-to-run spread of
\result{cifar10.noise.floor.lo}--\result{cifar10.noise.floor.hi} accuracy points
quoted in Sections~\ref{sec:experiments} and \ref{sec:discussion} was estimated
from repeated runs of a fixed configuration on this dataset, and is used only to
set a floor below which a single-cell difference should not be interpreted. It is
not a per-cell error bar: the headline cells in Table~\ref{tab:seeds} carry their
own measured standard deviations over 3 seeds, and those should be preferred
wherever both are available.
